\documentclass[10pt,a4paper]{amsart}
\usepackage[margin=19mm]{geometry}
\usepackage{amsmath,amssymb,mathtools}
\usepackage[scr=rsfs,cal=euler]{mathalfa}
\usepackage{xfrac}
\usepackage[unicode,hidelinks,bookmarks=false]{hyperref}
\newcommand{\ie}{i.e.\ }
\newcommand{\cf}{cf.\ }
\newcommand{\resp}{respectively\ }
\newcommand{\Subsection}{\subsection}
\providecommand{\nobreakdash}{\nobreak\hbox{-}}
\setlength{\emergencystretch}{2em}
\multlinegap0pt
\usepackage{bm}
\usepackage{bbm}
\usepackage{dsfont}
\usepackage{xspace}
\usepackage{cancel}
\usepackage{mathtools}
\usepackage{xcolor}
\usepackage{amsthm}
\makeatletter
\@ifundefined{thm}{\newtheorem{thm}{Thm}}{}
\@ifundefined{prop}{\newtheorem{prop}[thm]{Proposition}}{}
\makeatother
\title[Horizontal decompositions, II]{Horizontal decompositions, II}
\author{Paolo Lisca, Andrea Parma}
\date{Source: arXiv:2302.14606v2 (CC BY 4.0)}
\begin{document}
\maketitle

\noindent\emph{Source and licence.} Horizontal decompositions, II. Version arXiv:2302.14606v2 (CC BY 4.0), distributed under the Creative Commons Attribution 4.0 International licence, as recorded in the arXiv licence metadata for this exact version and shown on the abstract page https://arxiv.org/abs/2302.14606v2. Full rights notice: licenses/arxiv-2302.14606-CC-BY-4.0.txt.\par\smallskip

\noindent\textit{Editorial scope.} Counted unit: the Prop whose source label is \texttt{p:negcohn} in the article (source0.tex lines 862--866, source label \texttt{p:negcohn}), delivered together with the article's own complete original proof of it (source0.tex lines 868--933, one \texttt{proof} environment of 66 lines). No mathematics is written, repaired or re-derived: statement and proof are the article's own text line for line. Required context retained uncounted: e:negcohn (lines 857--859). Every definition, display and label of those lines is retained unchanged. Omitted from this package and not redistributed: 1--856, 860--861, 867--867, 934--1759. Nothing inside a retained excerpt is omitted or altered: every retained body is delivered exactly as the article's lines are written, so no omission receipt of any kind accompanies this package. Citations: the retained excerpts carry no citation command at all, so no bibliography entry of the article is transcribed and no reference is redistributed. Reference adaptation: none was needed and none was made; every reference inside the delivered unit resolves inside it, so no unresolved reference or citation is produced. Printed slips kept literal: no printed word, symbol, hypothesis or number of the counted statement or of its proof was altered. Layout and class: the article's own class is \texttt{amsart} with packages including geometry, graphicx, epstopdf, bm, soul, url, xcolor, pinlabel, biblatex, fourier, microtype, caption, verbatim; the wrapper is the lane's pinned \texttt{amsart} editorial wrapper at 10pt on A4, which restates the theorem-like environments the retained text needs and adds the article's own macro definitions that the retained text needs (none), with extra wrapper packages (bm, bbm, dsfont, xspace, cancel, mathtools, xcolor). The delivered numbering is the unit's own. Assets: no retained span contains a figure environment, a graphics-inclusion command, a PDF-page inclusion, a table or a TikZ picture; no image file is copied into this package, the package declares no asset dependency and no image file is redistributed with it. Licence: version arXiv:2302.14606v2 of this article is distributed under the Creative Commons Attribution 4.0 International licence, as recorded in the arXiv licence metadata for this exact version and shown on the abstract page https://arxiv.org/abs/2302.14606v2. A full rights notice is delivered at licenses/arxiv-2302.14606-CC-BY-4.0.txt. Characters: every retained excerpt is pure ASCII, so the delivered engine, XeLaTeX with its default Latin Modern text font, reports no missing character.
\begin{equation}\label{e:negcohn}
x^2 + xyz -a = y^2 + z^2,
\end{equation}

\begin{prop}\label{p:negcohn}
	Let $(x,y,z)$ be a solution of Equation~\eqref{e:negcohn} with $|x|$, $|y|$, $|z|$ pairwise distinct.  
	Then, $(x,y,z)$ is not weakly minimal. Indeed, mutation at the greatest element in absolute value makes the quantity $|x|+|y|+|z|$ strictly smaller. Moreover, if $\min(|x|,|y|,|z|)>1$ then mutation at the other 
	two elements makes the quantity $|x|+|y|+|z|$ strictly bigger. 
\end{prop}

\begin{proof} 
Note that if $(x,y,z)$ is a solution then so is $(x,z,y)$. Therefore, 
up to swapping $y$ and $z$ we may assume $|y|>|z|$. Moreover, it is easy to check that $xyz\neq 0$, 
because otherwise at least two among $x,y,z$ would be zero. We split the proof into the three cases
$|x| > |y| > |z|$, $|y| > |x| > |z|$ and $|y| > |z| > |x|$.

\smallskip\noindent
{\bf First case:} $|x| > |y| > |z|$. We claim that $|\hat x|<|x|$. Up to sign changes it suffices to assume $x,y>0$. This implies $z<0$, because $z>0$ gives the contradiction
\[
y^2+z^2 = x^2+xyz-a > y^2 + 2xz-a \ge y^2+2(z+2)z-a> y^2+z^2.
\]
We have $x(yz+x)=y^2+z^2+a>0$, therefore $\hat x = -yz - x < 0$, which implies, since $-yz>0$ and $-x<0$, $|\hat x|<|x|$. 
This shows that mutation at the greatest element in absolute value makes the quantity $|x|+|y|+|z|$ strictly smaller, in particular that $(x,y,z)$ is not weakly minimal. 
To prove the last part of the statement suppose $|z|>1$. 
Then, $y>2$ and  
\begin{equation}\label{e:zhatineq}
2|z| < xy < xy-z + |z| = \hat z + |z|,
\end{equation}
therefore $|z| < \hat z$.
\color{black}
Similarly, we have 
\[
2 x < |xz| \leq |xz - y| + |y| < |\hat y| + x,
\]
therefore $y < x < |\hat y|$. This concludes the proof of the statement in the first case. 

Before tackling the other cases we observe that up to sign changes we may 
assume $y>0$. Moreover, if $xz<0$ we have  
\[
y^2 -xyz + z^2 > x^2 -xyz + z^2 > x^2 \ge x^2-a,
\]
which is a contradiction. Hence, after a further sign change we may assume 
$x>0$ and $z>0$ as well. Therefore, from now on we assume $x,y,z>0$. 

\smallskip\noindent
{\bf Second case:} $|y|>|x|>|z|$. 
Note that $y\hat y = z^2-x^2+a \geq 0$ forces $x=2$, $z=1$ and $a=4$, in which case   
$y\hat y = 1$, which is impossible because $y>z\geq 1$. Therefore we have 
$y\hat y<0$, which implies $\hat y = xz-y<0$. 
It follows that $0 < xz < y$, hence $|\hat y| < |y|$. 
To prove the last part of the statement suppose $|z|>1$. Then, 
$|z|<|\hat z|$ follows as before from~\eqref{e:zhatineq}. 
Finally, we have 
$$
2|y| \leq |-yz - x| + |x| < |\hat x| + |y|,
$$	
hence $|x|<|y|<|\hat x|$. This proves the statement in the second case. 

	\smallskip\noindent
	{\bf Third case:} $|y|>|z|>|x|$.
Since $y\hat y = z^2-x^2+a>0$ we have $\hat y = xz-y>0$. If $\hat y \ge y$ then $xz \ge 2y$ 
and, using that $z\geq x+1\geq 2$, we get the contradictory inequalities
\[
y^2 + z^2 = x^2 + xyz-a > 2y^2-a \geq y^2 + (|z|+1)^2 - a > y^2 + z^2.
\]
Therefore $|y|=y > \hat y = |\hat y|$. 
To prove the last part of the statement observe that 
\[
2|x| < |yz| \leq |-yz-x|+|x| = |\hat x| + |x|, 
\]
hence $|x|<|\hat x|$. Using the extra assumption $|x|\geq 2$ we have  
$$
2|y| \leq |xy - z| + |z| < |\hat z| + |y|,
$$
hence $|z|<|y|<|\hat z|$. This proves the statement in the third case and concludes the proof of the proposition. 
\end{proof}

\end{document}
